% 全概率公式与Bayes公式的树状图。教学数据：
% 产线占比 P(B1)=.5,P(B2)=.3,P(B3)=.2；次品率 .01,.02,.04；
% 次品路径概率 .005,.006,.008，合计 P(D)=.019；P(B3|D)=.008/.019≈.421。
\begin{tikzpicture}[line width=.6pt,font=\fontsize{8}{10}\selectfont,>=Stealth,
  node distance=4mm,
  st/.style={draw=FigureBlue,fill=FigureBlueFill,rounded corners=1.8pt,inner sep=2.5pt,minimum width=9mm},
  df/.style={draw=FigureRed,fill=FigureRedFill,rounded corners=1.8pt,inner sep=2.5pt,minimum width=9mm},
  ok/.style={draw=FigureStroke,fill=white,rounded corners=1.8pt,inner sep=2.5pt,minimum width=9mm},
  pr/.style={fill=white,inner sep=1pt,font=\fontsize{7.5}{9}\selectfont}]
\node[st] (r) at (0,0) {产品};
\foreach \k/\y/\p/\q/\pq in {1/12/0.5/0.01/0.005,2/0/0.3/0.02/0.006,3/-12/0.2/0.04/0.008}{
  \node[st] (b\k) at (24mm,\y mm) {$B_{\k}$};
  \draw[FigureStroke] (r.east) -- node[pr,pos=.55]{$\p$} (b\k.west);
  \node[df] (d\k) at (50mm,\y mm+2.6mm) {次品};
  \node[ok] (o\k) at (50mm,\y mm-2.6mm) {合格};
  \draw[FigureRed] (b\k.east) -- node[pr,pos=.5,above=-.2mm]{$\q$} (d\k.west);
  \draw[FigureStroke] (b\k.east) -- (o\k.west);
  \node[FigureRed,anchor=west] at (61mm,\y mm+2.6mm) {$\p\times\q=\pq$};
}
\draw[decorate,decoration={brace,amplitude=3pt,raise=1pt},FigureRed]
  (88mm,16mm) -- (88mm,-10.5mm);
\node[FigureRed,anchor=west,align=left] at (91mm,2.8mm) {$\mathbb P(D)=0.019$};
\node[FigureInk,anchor=west,align=left] at (91mm,-3.2mm) {$\mathbb P(B_3\mid D)=\dfrac{0.008}{0.019}$};
\end{tikzpicture}
